% -------------------------------------------------------------------------
% WBS ID: RES-STR-LOD
% Deliverable: Structural Load Cases and Load Transfer
% Status: In progress
% Evidence status: Hydrodynamic load-transfer requirements established
% Review status: Not reviewed
% Last updated: 2026-07-19
% Dependencies: RES-PHY-LOD, RES-GEO-REP, RES-NUM-SPEC
% Downstream interfaces: RES-STR-FSI, RES-STR-SIM, RES-VER-MET, Software output schema
% -------------------------------------------------------------------------

\section{RES-STR-LOD --- Structural Load Cases and Load Transfer}
\label{sec:res-str-lod}

\subsection{Purpose and Scope}
\label{sec:res-str-lod-purpose}

This Deliverable defines the hydrodynamic load histories that must be
preserved for later structural response, damage or FSI analyses. It
does not activate material damage, deformable barriers, scour or
two-way fluid--structure coupling in the baseline tsunami-barrier
comparison.

\subsection{Research Questions}
\label{sec:res-str-lod-research-questions}

\begin{enumerate}
  \item Which pressure, shear, force, impulse and moment outputs are
    required from the local fluid solver?
  \item How are distributed fluid tractions reduced or mapped without
    losing resultant force and moment?
  \item Which structural quantities are stretch consumers rather than
    active baseline outputs?
\end{enumerate}

\subsection{Terminology and Definitions}
\label{sec:res-str-lod-definitions}

% TODO: Define the technical terminology, symbols, quantities and
% classifications used in this Deliverable.

\subsection{Literature Evidence}
\label{sec:res-str-lod-literature-evidence}

\textcite{baragamage2024fullycoupled} demonstrates the sequence from
three-dimensional tsunami fluid loading to structural motion and
feedback. Source role: supports preserving distributed load histories
for later FSI-capable analysis. Limitation: its reduced bridge
structure does not justify activating full deformable barrier analysis
in the current baseline.

\textcite{istrati2019trappedair} shows that pressure distribution can
alter component demand and moment even where total force is known.
Source role: supports retaining spatial pressure/traction fields, not
only resultant force. Project conclusion: the local hydrodynamic model
shall export complete load histories before any structural extension is
attempted.

\subsection{Governing Relations and Quantitative Definitions}
\label{sec:res-str-lod-governing-relations}

% TODO: Record relevant equations, dimensionless groups, empirical models,
% extraction rules or quantitative definitions.
%
% Use align environments for displayed equations.
% Every displayed equation must have a unique \label{eq:...}.
% Do not place punctuation after displayed equations.

\subsection{Material or Structural System}
\label{sec:res-str-lod-material-structural-system}

% TODO: Complete this RES-STR Domain-specific subsection.

\subsection{Load Cases}
\label{sec:res-str-lod-load-cases}

% TODO: Complete this RES-STR Domain-specific subsection.

\subsection{Hydrodynamic Traction Transfer}
\label{subsec:res_str_lod_hydrodynamic_traction_transfer}

The local three-dimensional fluid model supplies a transient traction field on the wetted structural surface. The traction exerted by the fluid on the solid is

\begin{align}
\bm{t}_{f\rightarrow s}
&=
\bm{\sigma}_{f}\bm{n}_{s}
\label{eq:res_str_lod_fluid_traction}
\end{align}

where \(\bm{n}_{s}\) is the outward unit normal of the current structural surface and the fluid Cauchy stress is

\begin{align}
\bm{\sigma}_{f}
&=
-p\bm{I}+\bm{\tau}_{f}
\label{eq:res_str_lod_fluid_cauchy_stress}
\end{align}

For a finite-element structural discretisation, the distributed interface traction is converted into the structural load vector through

\begin{align}
\bm{f}_{\mathrm{hyd}}(t)
&=
\int_{\Gamma_{fs}(t)}
\bm{N}_{s}^{\mathsf{T}}
\bm{t}_{f\rightarrow s}(t)
\,\mathrm{d}\Gamma
\label{eq:res_str_lod_hydrodynamic_load_vector}
\end{align}

The integration is performed over the current wetted interface \(\Gamma_{fs}(t)\). The transferred field must preserve the resultant force and moment and should preserve interface work when the fluid and structural surface discretisations do not match.

The tsunami FSI model of Baragamage and Wu demonstrates the sequence from three-dimensional fluid loading to structural motion and subsequent feedback to the fluid solver, although its structural representation is reduced rather than a full three-dimensional deformable continuum \autocite{baragamage2024fullycoupled}

\textbf{Selected approach.} Hydrodynamic loading is represented as a
distributed pressure and viscous traction field on the fixed barrier
surface during the active baseline. Reduction to a single resultant
force is not sufficient for later structural or damage studies.
Decision role: \cref{eq:res_str_lod_fluid_traction}--\cref{eq:res_str_lod_hydrodynamic_load_vector}
define the load-transfer contract. Limitation: the equations are a
handoff requirement and do not claim that a structural solver or FSI
run has been executed.

\subsection{Constitutive Behaviour}
\label{sec:res-str-lod-constitutive-behaviour}

% TODO: Complete this RES-STR Domain-specific subsection.

\subsection{Response and Stability Measures}
\label{sec:res-str-lod-response-stability-measures}

% TODO: Complete this RES-STR Domain-specific subsection.

\subsection{Damage and Failure Modes}
\label{sec:res-str-lod-damage-failure-modes}

% TODO: Complete this RES-STR Domain-specific subsection.

\subsection{Fluid--Structure Coupling Requirements}
\label{sec:res-str-lod-fluid-structure-coupling-requirements}

% TODO: Complete this RES-STR Domain-specific subsection.

\subsection{Candidate Approaches}
\label{sec:res-str-lod-candidate-approaches}

% TODO: Define the credible candidate approaches considered.

\subsection{Critical Comparison}
\label{sec:res-str-lod-critical-comparison}

% TODO: Compare candidates using physically and project-relevant criteria.
% Do not select an approach solely on popularity or implementation ease.

\subsection{Assumptions and Applicability}
\label{sec:res-str-lod-assumptions}

% TODO: State assumptions and the conditions under which they are valid.

\subsection{Limitations and Failure Conditions}
\label{sec:res-str-lod-limitations}

% TODO: State where the evidence, model or selected approach ceases to be
% reliable or applicable.

\subsection{Project-Specific Conclusions}
\label{sec:res-str-lod-conclusions}

The Studentship shall first stabilise hydrodynamic load extraction:
pressure, shear, force, moment and impulse histories on fixed barrier
patches. Material deformation and damage remain later consumers of
these histories.

\subsection{Selected Approach and Rationale}
\label{sec:res-str-lod-selected-approach}

The selected near-term approach is one-way hydrodynamic load export
from the local URANS--VOF model to structured output. One-way
CFD-to-FEA and two-way FSI are retained as stretch paths, not baseline
deliverables.

\subsection{Unresolved Decisions}
\label{sec:res-str-lod-unresolved-decisions}

Open decisions are load-mapping tolerances, time-history sampling
rate, structural mesh target, sign convention for moments and the
activation criteria for one-way structural analysis.

\subsection{Requirements and Handoffs}
\label{sec:res-str-lod-requirements}

Software shall emit \texttt{patch\_id}-resolved pressure, shear,
resultant force, impulse, overturning moment, torsional moment,
centre-of-pressure and integration-window metadata. RES-VER-MET shall
define tolerances for force/moment consistency before RES-STR-SIM or
RES-STR-FSI consumes the loads.

\subsection{Deliverable Completion Record}
\label{sec:res-str-lod-completion-record}

% TODO: Record whether the Deliverable is:
%
% - Not started
% - In progress
% - Draft complete
% - Reviewed
% - Accepted
%
% Acceptance requires:
%
% - the research questions to be answered;
% - claims to be supported by appropriate evidence;
% - assumptions and limitations to be explicit;
% - the project-specific conclusion to be stated;
% - downstream requirements to be traceable;
% - unresolved decisions to be closed or formally deferred.
